% =====================================================================
% Homework 03 --- Parametrizacao e geracao de perfis
% Incluido por main.tex. Tabelas em tex_airfoil_mod/ e figuras em
% resultados_airfoil_mod/ (mesmo nivel que este arquivo e o main.tex).
% Codigos em airfoil_mod/.
% =====================================================================

\section{Ponto de projeto da seção de referência}
\label{sec:ponto-projeto}

O primeiro passo do laboratório é fixar o ponto de projeto em que o
aerofólio será otimizado. A aeronave de referência é a NJ-0502 com os
inputs já atualizados para o ótimo do Lab~02. Adotou-se o
\textbf{meio do cruzeiro}: altitude e Mach de cruzeiro da missão de
projeto, e o peso médio aritmético entre o início e o fim dessa etapa da missão.

\subsection{Altitude, Mach e área de referência}

Na missão de projeto da NJ-0502,
\begin{equation}
h = 35\,000~\mathrm{ft},
\qquad
M_{\infty} = 0{,}85,
\qquad
S_{\mathrm{ref}} = S_w = 386{,}48~\mathrm{m}^2.
\end{equation}
A área de referência é a área alar otimizada no Lab~02.

\subsection{Atmosfera no ponto de projeto}

Com a atmosfera padrão ISA em $h = 35\,000$~ft ($10\,668$~m),
\begin{equation}
\rho_{\infty} = 0{,}3805~\mathrm{kg/m}^3,
\qquad
a_{\infty} = 296{,}59~\mathrm{m/s},
\qquad
\mu_{\infty} = 1{,}445\times 10^{-5}~\mathrm{Pa\,s}.
\end{equation}
A velocidade de voo correspondente é
\begin{equation}
V_{\infty} = M_{\infty}\,a_{\infty}
= 0{,}85\times 296{,}59
= 252{,}10~\mathrm{m/s}.
\end{equation}

\subsection{Peso no meio do cruzeiro}

O \emph{Design Tools} decompõe a missão em frações de massa. Para a
categoria \emph{transport},
\begin{equation}
\begin{aligned}
M_{f,\mathrm{start}} &= 0{,}990,&
M_{f,\mathrm{taxi}} &= 0{,}990,&
M_{f,\mathrm{TO}} &= 0{,}995,&
M_{f,\mathrm{climb}} &= 0{,}980.
\end{aligned}
\end{equation}
Com o MTOW convergido $W_0 = 276\,276$~kgf, o peso no \emph{início} do
cruzeiro é
\begin{equation}
W_{\mathrm{ini}}
= W_0\,M_{f,\mathrm{start}}\,M_{f,\mathrm{taxi}}\,M_{f,\mathrm{TO}}\,M_{f,\mathrm{climb}}
= 264\,036~\mathrm{kgf}.
\end{equation}
A fração de cruzeiro vem da equação de Breguet implementada em
\code{fuel\_weight},
\begin{equation}
M_{f,\mathrm{cruise}} = 0{,}6992,
\end{equation}
de modo que o peso no \emph{fim} do cruzeiro fica
\begin{equation}
W_{\mathrm{fim}} = W_{\mathrm{ini}}\,M_{f,\mathrm{cruise}}
= 184\,607~\mathrm{kgf}.
\end{equation}
O ponto de projeto adotado é a média aritmética entre esses extremos:
\begin{equation}
W
= \frac{W_{\mathrm{ini}} + W_{\mathrm{fim}}}{2}
= 224\,321~\mathrm{kgf}.
\end{equation}

\subsection{Coeficiente de sustentação da aeronave}

Em voo nivelado,
\begin{equation}
C_L
= \frac{2W}{\rho_{\infty}\,S_{\mathrm{ref}}\,V_{\infty}^{2}}
= 0{,}4710.
\end{equation}
Seguindo o enunciado, a seção de referência herda esse valor:
\begin{equation}
c_{l\mathrm{ref}} = C_L = 0{,}4710.
\end{equation}

\subsection{Corda e espessura da seção de referência}

A corda de referência é a corda média aerodinâmica (MAC) da asa
otimizada,
\begin{equation}
c_{\mathrm{ref}} = c_{m,w}
= \frac{2}{3}\,c_{r,w}\,\frac{1+\lambda_w+\lambda_w^{2}}{1+\lambda_w}
= 6{,}918~\mathrm{m},
\end{equation}
com $\lambda_w = 0{,}20$. A espessura relativa mínima para a otimização
é a média entre raiz e ponta da asa da NJ-0502,
\begin{equation}
\left(\frac{t}{c}\right)_{\mathrm{ref}}
= \frac{(t/c)_{r} + (t/c)_{t}}{2}
= \frac{0{,}16 + 0{,}08}{2}
= 0{,}12.
\end{equation}

\subsection{Reynolds da seção de referência}

\begin{equation}
Re_{\mathrm{ref}}
= \frac{\rho_{\infty}\,V_{\infty}\,c_{\mathrm{ref}}}{\mu_{\infty}}
= 4{,}592\times 10^{7}.
\end{equation}

\subsection{Tabela-resumo}

Os dados da seção de referência estão reunidos na Tabela~\ref{tab:ponto}.

\begin{table}[H]
\centering
\caption{Dados para a seção de referência.}
\label{tab:ponto}
\small
\begin{tabular}{lll}
\toprule
Parâmetro & Valor & Descrição \\
\midrule
$h$ & $35\,000$ & Altitude do ponto de projeto [ft] \\
$M_{\infty}$ & $0{,}85$ & Mach do ponto de projeto \\
$W$ & $224\,321{,}2$ & Peso da aeronave no ponto de projeto [kgf] \\
$S_{\mathrm{ref}}$ & $386{,}48$ & Área de referência da aeronave [m$^2$] \\
$\rho_{\infty}$ & $0{,}3805$ & Densidade do ar no ponto de projeto [kg/m$^3$] \\
$a_{\infty}$ & $296{,}59$ & Velocidade do som no ponto de projeto [m/s] \\
$c_{\mathrm{ref}}$ & $6{,}918$ & Corda da seção de referência [m] \\
$c_{l\mathrm{ref}}$ & $0{,}4710$ & Coeficiente de sustentação da seção de referência \\
$Re_{\mathrm{ref}}$ & $4{,}592 \times 10^{7}$ & Reynolds para a seção de referência \\
$(t/c)_{\mathrm{ref}}$ & $0{,}12$ & Espessura relativa base da seção de referência \\
\bottomrule
\end{tabular}

\end{table}

% ---------------------------------------------------------------------
\section{Parametrização e Geração de Perfis}
\label{sec:airfoil-mod}
% ---------------------------------------------------------------------

\subsection{Objetivo}

Explique aqui qual parametrização foi usada no laboratório
(por exemplo, NACA, CST ou outra), quais variáveis geométricas foram
adotadas e qual era o objetivo desta etapa.

\subsection{Organização dos códigos}

Os códigos desta seção estão em
\code{designTool/lab03\_perfil/airfoil\_mod/}:
\begin{itemize}
\item \code{run\_1\_ponto\_projeto.py} --- ponto de projeto e Tabela~\ref{tab:ponto};
\item demais rotinas de geração/ajuste de perfil.
\end{itemize}

\subsection{Metodologia}

Registre a definição do perfil, os parâmetros geométricos escolhidos,
o procedimento de geração das coordenadas e, se houver, o método de ajuste
do perfil de referência.

\subsection{Resultados}

Insira aqui as figuras e tabelas desta etapa.

\begin{figure}[H]
\centering
% \includegraphics[width=0.85\textwidth]{resultados_airfoil_mod/exemplo.png}
\caption{Substitua esta legenda por uma descrição da figura da etapa de
parametrização.}
\label{fig:airfoil-mod-exemplo}
\end{figure}

\subsection{Discussão}

Discuta a qualidade do ajuste, a interpretação dos parâmetros e eventuais
vantagens ou limitações da parametrização usada.
